\section{Getting Started}
%%%%%%%%%%%%%%%%%%%%%%%%%%%%%%%%%%%%%%%%%%%%%%%%%%%%%%%%%%%%%%%%%%%%%%%%%%%%%%%%%%
\begin{frame}[fragile]\frametitle{Installation: All Platforms}
  \begin{itemize}
    \item \textbf{Quick Install (macOS / Linux):}
    \lstinline|curl -fsSL https://claude.ai/install.sh | bash|
    \item \textbf{macOS (Homebrew):}
    \lstinline|brew install --cask claude-code|
    \item \textbf{Windows (PowerShell):}
    \lstinline|irm https://claude.ai/install.ps1 | iex|
    \item \textbf{Windows (CMD):}
    \lstinline|curl -fsSL https://claude.ai/install.cmd -o install.cmd && install.cmd|
    \item \textbf{npm (legacy, deprecated since v2.1.15):}
    \lstinline|npm install -g @anthropic-ai/claude-code|
    \item Migrate from npm to native binary: \lstinline|claude install|
    \item Verify installation: \lstinline|claude doctor|
    \item Key file locations (macOS/Linux):
      \begin{itemize}
        \item Binary: \lstinline|~/.local/bin/claude|
        \item Global config: \lstinline|~/.claude/settings.json|
        \item Project config: \lstinline|.claude/settings.json| (repo root)
        \item Auth/state: \lstinline|~/.claude.json|
        \item Sessions: \lstinline|~/.claude/sessions/|
      \end{itemize}
  \end{itemize}
\end{frame}

%%%%%%%%%%%%%%%%%%%%%%%%%%%%%%%%%%%%%%%%%%%%%%%%%%%%%%%%%%%%%%%%%%%%%%%%%%%%%%%%%%
\begin{frame}[fragile]\frametitle{Getting Started: Authentication}

  \begin{itemize}
    \item \textbf{Launch:} Run \texttt{claude} in any project directory. The REPL loads; your first step is authentication.
    \item \textbf{Authenticate:} Run \texttt{claude auth login} → browser opens → sign in with Anthropic account.
    \item \textbf{Check status:} \texttt{claude auth status} shows current account, plan, and token expiry.
    \item \textbf{Config format:} JSON-based \texttt{settings.json}; supports project-level and global scopes.
    \item \textbf{First Step After Auth:} Run \texttt{/init} to ground Claude in your project --- scans the codebase and creates \texttt{CLAUDE.md}.
  \end{itemize}

\end{frame}

%%%%%%%%%%%%%%%%%%%%%%%%%%%%%%%%%%%%%%%%%%%%%%%%%%%%%%%%%%%%%%%%%%%%%%%%%%%%%%%%%%
\begin{frame}[fragile]\frametitle{Access Options in Claude Code}

\textbf{Claude Code supports three access paths.}

\vspace{0.3cm}
\textbf{Claude Pro / Max Subscription}

\begin{itemize}
    \item Single monthly plan covers both claude.ai web and CLI usage
    \item Predictable cost; no per-token billing surprises
    \item Includes rate-limited access to all Claude models
\end{itemize}

\vspace{0.2cm}
\textbf{Anthropic Console (API Billing)}

\begin{itemize}
    \item Pay per token; no subscription; full API access
    \item Claude Code workspace auto-created at platform.claude.com
    \item Pricing (per 1M tokens, Sonnet 4.6): \$3.00 input / \$15.00 output
    \item Haiku 4.5: \$0.80 input / \$4.00 output (ideal for subagents)
\end{itemize}

\vspace{0.2cm}
\textbf{Enterprise Platforms}
\begin{itemize}
    \item AWS Bedrock, Google Vertex AI, Microsoft Foundry
    \item Integrates with existing cloud billing and compliance controls
    \item Data residency options; no data retention by default
\end{itemize}

\end{frame}

%%%%%%%%%%%%%%%%%%%%%%%%%%%%%%%%%%%%%%%%%%%%%%%%%%%%%%%%%%%%%%%%%%%%%%%%%%%%%%%%%%
\begin{frame}[fragile]\frametitle{Enterprise Auth: Bedrock, Vertex, Foundry}

\textbf{Benefits for enterprise teams:}
\begin{itemize}
    \item Leverage existing cloud spend commitments and quotas
    \item Data stays within your cloud region (compliance)
    \item Works behind firewalls with proxy support
    \item Entra ID (Azure) or IAM (AWS) for identity management
\end{itemize}

\textbf{Claude Code can route through your existing cloud infrastructure:}

\begin{lstlisting}
# AWS Bedrock
export CLAUDE_CODE_USE_BEDROCK=1
export AWS_REGION=us-east-1
export AWS_PROFILE=your-profile

# Google Vertex AI
export CLAUDE_CODE_USE_VERTEX=1
export CLOUD_ML_REGION=us-east5
export ANTHROPIC_VERTEX_PROJECT_ID=your-project

# Microsoft Foundry (Azure)
export CLAUDE_CODE_USE_FOUNDRY=1
export ANTHROPIC_FOUNDRY_RESOURCE=your-resource-name

# Corporate proxy
export HTTPS_PROXY='https://proxy.example.com:8080'
\end{lstlisting}

\end{frame}

%%%%%%%%%%%%%%%%%%%%%%%%%%%%%%%%%%%%%%%%%%%%%%%%%%%%%%%%%%%%%%%%%%%%%%%%%%%%%%%%%%
\begin{frame}[fragile]\frametitle{Plan vs Normal vs Auto-Accept Mode}

Switch between them using \texttt{Shift+Tab} (cycles through all three).

\textbf{Plan Mode:}
  \begin{itemize}
    \item Read-only --- cannot make changes to your code
    \item Analyzes and explores your codebase with full reasoning
    \item Writes a plan to \texttt{.claude/plans/} for your review
    \item Available tools: Read, Glob, Grep, WebSearch, WebFetch
    \item Perfect for: architecture decisions, complex refactors, unfamiliar codebases
  \end{itemize}

\textbf{Normal Mode (Default):}
  \begin{itemize}
    \item Full access --- can read, write, and modify files
    \item Prompts you to approve each file change
    \item Best for: careful development where you want per-edit review
  \end{itemize}

\textbf{Auto-Accept Mode:}
  \begin{itemize}
    \item All edits applied without per-file confirmation
    \item Claude shows a running diff and cost in the status bar
    \item Interrupt at any time with \texttt{Ctrl+C}
    \item Best for: bulk scaffolding and trusted, well-specified tasks
  \end{itemize}

\textbf{Pro Tip:} Always start complex features in Plan mode to review the approach, then switch to Normal or Auto-Accept for implementation.
\end{frame}

%%%%%%%%%%%%%%%%%%%%%%%%%%%%%%%%%%%%%%%%%%%%%%%%%%%%%%%%%%%%%%%%%%%%%%%%%%%%%%%%%%
\begin{frame}[fragile]\frametitle{Context Mentions \& Image Support}
  \begin{itemize}
    \item \textbf{\texttt{@file:path}} --- focus the agent on a specific file:
      \lstinline|@file:main.py fix the bug where user input isn't sanitized|
    \item \textbf{\texttt{@file:}} without path --- opens fuzzy file picker for the project.
    \item \textbf{Onboarding power:} Instead of manually \texttt{grep}-ing and reading files, just ask. Especially useful when exploring an unfamiliar codebase.
    \item \textbf{Image paste/drag-and-drop:} Paste screenshots or mockups directly into the terminal (\texttt{Cmd+V} on macOS). Claude sees the image and can implement UI from a design mockup, diagnose a browser error, or read a diagram.
    \item \textbf{Shift+Tab to Plan mode:} For complex changes, switch to Plan mode, describe the feature, and review the proposed approach before any code is written.
    \item \textbf{Alt+T for Extended Thinking:} Toggle deep reasoning mode for hard architectural or algorithmic problems.
  \end{itemize}

\end{frame}

%%%%%%%%%%%%%%%%%%%%%%%%%%%%%%%%%%%%%%%%%%%%%%%%%%%%%%%%%%%%%%%%%%%%%%%%%%%%%%%%%%
\begin{frame}[fragile]\frametitle{Session Management}
  \begin{itemize}
    \item \textbf{\texttt{claude -c}} (or \texttt{--continue}) --- resumes the most recent session with full context.
    \item \textbf{\texttt{claude --resume}} --- opens a session picker; select by name or number.
    \item \textbf{Name sessions:} Use \texttt{/rename auth-refactor} inside a session, or \texttt{claude -n "auth-refactor"} at launch for easy retrieval.
    \item \textbf{Fork a session:} \texttt{claude --resume <name> --fork-session} --- branch to try an alternative approach without losing the original.
    \item \textbf{PR-linked sessions:} \texttt{claude --from-pr 42} --- start a session linked to a specific pull request. Sessions auto-link when you create a PR via \texttt{gh} during a session.
    \item \textbf{\texttt{/compact}} --- compresses the current session context to save tokens while preserving key decisions and code snippets.
    \item \textbf{\texttt{/cost}} --- shows token usage and running session cost at any time.
  \end{itemize}
{\tiny (Ref: Anthropic, ``Claude Code CLI Reference'', docs.claude.com/en/docs/claude-code)}
\end{frame}

%%%%%%%%%%%%%%%%%%%%%%%%%%%%%%%%%%%%%%%%%%%%%%%%%%%%%%%%%%%%%%%%%%%%%%%%%%%%%%%%%%
\begin{frame}[fragile]\frametitle{Non-Interactive Mode and CI/CD}
  \begin{itemize}
    \item \textbf{\texttt{claude -p}} (print mode) --- executes a single query and exits; ideal for scripting and CI.
    \item \textbf{Structured output:} \texttt{--output-format json} returns cost, duration, session ID, and result for parsing in pipelines.
    \item \textbf{Control agentic depth:} \texttt{--max-turns 10} prevents runaway loops; \texttt{--allowedTools} restricts what Claude can do.
    \item \textbf{GitHub Action:} Use \texttt{claude -p} in CI pipelines --- auto-review PRs, generate changelogs, run AI-assisted linting.
  \end{itemize}

\begin{lstlisting}
# One-shot task
claude -p "Add docstrings to all public functions in src/"

# Structured output for scripts
claude -p "Count lines by file type" --output-format json

# Continue last session non-interactively
claude -c -p "Now add unit tests for those functions"

# In GitHub Actions:
claude -p "Review changed files for bugs: $CHANGED.
  Post findings as a PR comment via GitHub MCP."
\end{lstlisting}

\end{frame}

\section{Use Cases \& Demo}
%%%%%%%%%%%%%%%%%%%%%%%%%%%%%%%%%%%%%%%%%%%%%%%%%%%%%%%%%%%
%% USE CASES
%%%%%%%%%%%%%%%%%%%%%%%%%%%%%%%%%%%%%%%%%%%%%%%%%%%%%%%%%%%
\begin{frame}[fragile]\frametitle{Real-World Use Cases}
  \begin{itemize}
    \item \textbf{Understanding a Codebase:}
    \lstinline|Explain the authentication flow in this project|
    Claude analyzes your files and explains how auth works with actual file references.

    \item \textbf{Adding a Feature:}
    \lstinline|Add a dark mode toggle to the settings page|
    Claude finds the settings page, identifies the theme system, implements the toggle, and updates styles.

    \item \textbf{Debugging:}
    \lstinline|The login form is not submitting. Error: [paste]|
    Claude traces the issue, checks relevant files, and suggests targeted fixes.

    \item \textbf{Multi-File Refactoring:} Rename, restructure, or migrate patterns across the entire codebase atomically.
    \item \textbf{Documentation Generation:} Auto-generate README, API docs, and inline comments from code context.
    \item \textbf{Test Generation:} Scaffold unit/integration tests with awareness of existing test patterns.
    \item \textbf{Security Auditing:} Identify unsanitized inputs, API key leaks, or injection vulnerabilities.
  \end{itemize}
\end{frame}


%%%%%%%%%%%%%%%%%%%%%%%%%%%%%%%%%%%%%%%%%%%%%%%%%%%%%%%%%%%%%%%%%%%%%%%%%%%%%%%%%%
\begin{frame}[fragile]\frametitle{}

\begin{center}
{\Large Demo}
\end{center}
\end{frame}

%%%%%%%%%%%%%%%%%%%%%%%%%%%%%%%%%%%%%%%%%%%%%%%%%%%%%%%%%%%%%%%%%%%%%%%%%%%%%%%%%%
\begin{frame}[fragile]\frametitle{Initialize \& Understand}
  \begin{itemize}
    \item \textbf{Goal:} Show how Claude Code scans and understands a new project from scratch.
    \item \textbf{Setup:} Open terminal inside any Git repository.
    \item \textbf{Launch:} \lstinline|$ claude|
    \item \textbf{Command:} \lstinline|> /init|
    \item \textbf{What Happens:} Claude scans all local files, infers the stack (languages, deps, structure), and writes \texttt{CLAUDE.md} --- the AI's instruction manual for this project.
    \item \textbf{Key Point:} Every subsequent interaction is grounded in this context --- no hallucinated file paths or invented conventions.
    \item \textbf{Auto Memory:} As Claude works, it appends learnings (build commands, debugging insights, discovered patterns) to \texttt{CLAUDE.md} automatically.
  \end{itemize}
\end{frame}

%%%%%%%%%%%%%%%%%%%%%%%%%%%%%%%%%%%%%%%%%%%%%%%%%%%%%%%%%%%%%%%%%%%%%%%%%%%%%%%%%%
\begin{frame}[fragile]\frametitle{Fix \& Refactor}
  \begin{itemize}
    \item \textbf{Goal:} Demonstrate AI-driven code modification on a real file.
    \item \textbf{Command:} \lstinline|> @file:main.py fix the bug where user input isn't sanitized|
    \item \textbf{Agent Action:} Reads \texttt{main.py}, identifies the unsanitized input vector, proposes a targeted patch.
    \item \textbf{Review Flow:}
      \begin{itemize}
        \item Press \texttt{Shift+Tab} → switch to \textbf{Plan} mode: inspect the proposed diff before anything touches disk.
        \item Press \texttt{Shift+Tab} → switch to \textbf{Normal} mode: approve each file change individually.
        \item Press \texttt{Shift+Tab} → switch to \textbf{Auto-Accept} mode: apply all changes automatically.
      \end{itemize}
    \item \textbf{Key Point:} You stay in control --- the agent proposes, you decide the approval level.
  \end{itemize}
\end{frame}

\section{Configuration}
%%%%%%%%%%%%%%%%%%%%%%%%%%%%%%%%%%%%%%%%%%%%%%%%%%%%%%%%%%%%%%%%%%%%%%%%%%%%%%%%%%
\begin{frame}[fragile]\frametitle{}

\begin{center}
{\Large Configuration}
\end{center}
\end{frame}


%%%%%%%%%%%%%%%%%%%%%%%%%%%%%%%%%%%%%%%%%%%%%%%%%%%%%%%%%%%%%%%%%%%%%%%%%%%%%%%%%%
\begin{frame}[fragile]\frametitle{Overview}
  \begin{itemize}
    \item Claude Code is fully configurable at \textbf{two primary scopes}: User (\texttt{\textasciitilde/.claude/}) and Project (\texttt{.claude/} in repo root).
    \item \textbf{Config priority} (highest → lowest): Enterprise managed → CLI flags → Local project \texttt{.claude/settings.local.json} → Shared project \texttt{.claude/settings.json} → User \texttt{\textasciitilde/.claude/settings.json}.
    \item \textbf{Subdirectories}: \texttt{agents/}, \texttt{skills/}, \texttt{commands/}, \texttt{rules/}, \texttt{plans/}.
	\item \textbf{Meanings}:
	  \begin{itemize}
			\item settings.json : Shared permission rules
			\item settings.local.json : Personal permissions (gitignored)
			\item rules  : Auto-enforced contextual rules
			\item commands  : Slash commands for common workflows
			\item agents  : Specialized subagent personas
			\item skills  :  Auto-triggered workflow checks
      \end{itemize}
    \item \textbf{Schema-validated:} Config validated against \texttt{json.schemastore.org/claude-code-settings.json}; editors autocomplete fields.
    \item \textbf{Team sharing:} \texttt{.claude/settings.json} and \texttt{.mcp.json} are committed to Git; \texttt{.claude/settings.local.json} is personal (add to \texttt{.gitignore}).
    \item Each section below covers: what ships built-in → what you can customize → how to create your own.
  \end{itemize}

\textbf{The core problem:} Context is a finite budget. Dumping everything into one file degrades quality --- the agent buries critical constraints under noise. Claude Code solves this with \textit{tiered, purposeful} context loading.
\end{frame}


%%%%%%%%%%%%%%%%%%%%%%%%%%%%%%%%%%%%%%%%%%%%%%%%%%%%%%%%%%%
%% Customization Layers
%%%%%%%%%%%%%%%%%%%%%%%%%%%%%%%%%%%%%%%%%%%%%%%%%%%%%%%%%%%

%%%%%%%%%%%%%%%%%%%%%%%%%%%%%%%%%%%%%%%%%%%%%%%%%%%%%%%%%%%%%%%%%%%%%%%%%%%%%%%%%%
\begin{frame}[fragile]\frametitle{Customization Layers: The Mental Model}
  \begin{itemize}
    \item \textbf{CLAUDE.md: ``The Onboarding Doc'':} \textit{Always-on}, repo-committed, team-shared. Tells \textit{any} agent how this project works: stack, build commands, test steps, conventions. Auto-updated by Claude as it learns. Think of it as the README for AI.
    \item \textbf{Rules: ``Standing Orders'':} \textit{Always-on} for matched paths, injected into every session. Placed in \texttt{.claude/rules/} with glob patterns. Encode constraints near high-risk directories (\texttt{src/auth/}, \texttt{infra/}).
    \item Claude Code has five extension mechanisms: commands, skills, subagents, hooks, and plugins. \textbf{Commands} = what to do. \textbf{Skills} = how to do it. \textbf{Subagents} = who does it. \textbf{Hooks} = guardrails that always run. \textbf{Plugins} = distributable bundles of all the above.
    \item \textbf{Skills: ``The On-Call Expert'':} \textit{Lazy-loaded}: Claude reads only the name and description at startup, fetches full content only when the task is relevant. Prevents token waste on irrelevant instructions.
    \item \textbf{Subagents: ``The Specialist Persona'':} Own model tier, tool set, and system prompt. Run in isolated contexts. Return only summaries to the parent --- preventing context bloat.
    \item \textbf{Hooks: ``The Guardrails'':} Shell/prompt/agent handlers that fire automatically on lifecycle events. They are \textit{deterministic} --- not probabilistic like prompts.
  \end{itemize}
\end{frame}

%%%%%%%%%%%%%%%%%%%%%%%%%%%%%%%%%%%%%%%%%%%%%%%%%%%%%%%%%%%%%%%%%%%%%%%%%%%%%%%%%%
\begin{frame}[fragile]\frametitle{Customization Layers: When to Use What}

\textbf{The layering principle:} Rules = what to always know. Skills = what to know when needed. Subagents = who does the work. Hooks = rules that always enforce.

  \begin{itemize}
    \item \textbf{Use CLAUDE.md when:} The instruction is \textit{project-wide}, stable, and useful to every agent, every session: build steps, test commands, folder conventions, PR rules. Commit to Git; share with the team. Auto-generate with \texttt{/init}.
    \item \textbf{Use Rules when:} The instruction applies to specific paths or is always relevant but not in the main \texttt{CLAUDE.md}: directory constraints, security invariants, package-level conventions. Global rules in \texttt{\textasciitilde/.claude/CLAUDE.md} stay personal.
    \item \textbf{Use Skills when:} The instruction is \textit{task-specific} and only occasionally needed: a \texttt{pytest-patterns} skill for test generation, a \texttt{refactor-patterns} skill for code cleanup. Keeps context lean.
    \item \textbf{Use Custom Subagents when:} You need a \textit{different set of capabilities}: a \texttt{review} subagent that can't write files, a \texttt{debug} subagent with only bash+read, a \texttt{docs} subagent using a premium model.
    \item \textbf{Use Commands (\texttt{/}) when:} You have a \textit{repetitive prompt}: a parameterized template you type often. Commands invoke subagents, pass arguments, and save keystrokes.
    \item \textbf{Use Hooks when:} You need \textit{deterministic} enforcement: auto-format after every edit, block dangerous commands, run tests before stopping.
  \end{itemize}
\end{frame}

%%%%%%%%%%%%%%%%%%%%%%%%%%%%%%%%%%%%%%%%%%%%%%%%%%%%%%%%%%%%%%%%%%%%%%%%%%%%%%%%%%
\begin{frame}[fragile]\frametitle{Customization Layers: Summary Table}

\begin{center}
\adjustbox{max width=\linewidth}{%
      \begin{tabular}{|l|l|l|l|}
        \hline
        \textbf{Layer} & \textbf{Loaded} & \textbf{Scope} & \textbf{Purpose} \\
        \hline
        CLAUDE.md & Always & Project / Personal & Project context \\
        Rules & Path-matched & Project / Personal & Constraints \\
        Skills & On-demand & Any & Domain expertise \\
        Subagents & When invoked & Session role & Specialist work \\
        Commands & When typed & Project / Global & Shortcuts \\
        Hooks & Automatically & Session & Deterministic rules \\
        Plugins & On install & Project / Global & Distribution \\
        \hline
      \end{tabular}
}
\end{center}

\end{frame}


%%%%%%%%%%%%%%%%%%%%%%%%%%%%%%%%%%%%%%%%%%%%%%%%%%%%%%%%%%%%%%%%%%%%%%%%%%%%%%%%%%
\begin{frame}[fragile]\frametitle{CLAUDE.md}
CLAUDE.md helps Claude Code understand: project structure, coding conventions, technology stack, build commands, and more.

\begin{lstlisting}
## Project: My App

## Tech Stack
- Next.js 14 with App Router
- TypeScript, Tailwind CSS

## Build Commands
- Install: `npm install`
- Dev:     `npm run dev`
- Tests:   `npm test`
- Lint:    `npm run lint`

## Conventions
- Use functional components with hooks
- Type hints on all exported functions
- No hardcoded secrets: always use env vars

## File Map
- /app       - Next.js app router pages
- /components - Reusable UI components
- /lib       - Shared utilities
\end{lstlisting}
\end{frame}

%%%%%%%%%%%%%%%%%%%%%%%%%%%%%%%%%%%%%%%%%%%%%%%%%%%%%%%%%%%%%%%%%%%%%%%%%%%%%%%%%%
\begin{frame}[fragile]\frametitle{Rules: Persistent AI Instructions}
  \begin{itemize}
    \item \textbf{What:} Markdown instructions injected into every LLM context matching the path pattern. Think of them as standing orders: coding style, constraints, security invariants.
    \item \textbf{Built-in bootstrap:} \texttt{/init} auto-generates \texttt{CLAUDE.md} by scanning your repo: languages, deps, folder structure, patterns.
    \item \textbf{Auto Memory:} Claude automatically appends learnings to \texttt{CLAUDE.md} during sessions. Use \texttt{/memory} to view or clear them.
    \item \textbf{Scope:}
      \begin{itemize}
        \item \textit{Project rules}: \texttt{./CLAUDE.md} --- commit to Git; shared across team.
        \item \textit{Global rules}: \texttt{\textasciitilde/.claude/CLAUDE.md} --- personal, not committed.
        \item \textit{Directory rules}: Nested \texttt{CLAUDE.md} files in subdirectories; loaded for work in that subtree.
        \item \textit{Rules files}: \texttt{.claude/rules/*.md} with \texttt{paths:} glob patterns for targeted loading.
      \end{itemize}
    \item \textbf{Create custom:} Write any Markdown file, reference it in \texttt{settings.json}. No special syntax: plain instructions work.
  \end{itemize}

\begin{lstlisting}
# .claude/rules/auth-rules.md
paths: ["src/auth/**", "tests/auth/**"]
NEVER store tokens in localStorage.
Always use httpOnly cookies for auth sessions.
Run the full auth test suite after any change here.
\end{lstlisting}
\end{frame}

%%%%%%%%%%%%%%%%%%%%%%%%%%%%%%%%%%%%%%%%%%%%%%%%%%%%%%%%%%%%%%%%%%%%%%%%%%%%%%%%%%
\begin{frame}[fragile]\frametitle{Subagents: Built-in \& Custom}
  \begin{itemize}
    \item \textbf{What:} Specialist agents that run in isolated contexts, do focused work, and return only summaries to the parent conversation.
    \item \textbf{Isolation benefit:} Subagent exploration does \textit{not} bloat your main context; only the conclusion returns.
    \item \textbf{Built-in subagent types (auto-selected by Claude):}
      \begin{itemize}
        \item \texttt{Explore}: read-only codebase navigator; fast file/pattern search.
        \item \texttt{Plan}: architect mode; designs approach without modifying files.
        \item \texttt{General}: full tool access; multi-step parallel tasks.
      \end{itemize}
    \item \textbf{Common custom subagents:} \texttt{review} (read-only + glob), \texttt{debug} (bash + read), \texttt{docs} (write, no shell), \texttt{specs} (write markdown only).
    \item \textbf{Create custom:} Write a Markdown file in \texttt{.claude/agents/<n>.md}. Claude auto-routes to a subagent when its \texttt{description} matches the task.
    \item \textbf{Subagent frontmatter} fields: \texttt{name}, \texttt{description}, \texttt{model}, \texttt{allowed-tools}.
    \item \textbf{Parallel execution:} Up to 10 subagents can run simultaneously for parallel research or review tasks.
  \end{itemize}
\end{frame}

%%%%%%%%%%%%%%%%%%%%%%%%%%%%%%%%%%%%%%%%%%%%%%%%%%%%%%%%%%%%%%%%%%%%%%%%%%%%%%%%%%
\begin{frame}[fragile]\frametitle{Built-in Tools}
  \begin{itemize}
    \item \textbf{What:} Functions Claude calls to act on your codebase. Core tools are enabled by default.
    \item \textbf{Core Built-in Tools:}
      \begin{itemize}
        \item \texttt{Read}: read file contents (supports line ranges)
        \item \texttt{Edit} / \texttt{Write}: exact string replacement; create/overwrite files
        \item \texttt{Bash}: execute shell commands (\texttt{git}, \texttt{npm}, \texttt{pytest}, etc.)
        \item \texttt{Glob}: find files by pattern (\texttt{**/*.py})
        \item \texttt{Grep}: regex search across codebase
        \item \texttt{WebFetch} / \texttt{WebSearch}: fetch URLs or search during sessions
        \item \texttt{TodoWrite}: manage task lists for multi-step work
      \end{itemize}
    \item \textbf{Restricting tools per subagent:} Use \texttt{allowed-tools} in subagent frontmatter to scope access. A \texttt{review} subagent with only \texttt{Read, Glob, Grep} cannot accidentally delete files.
    \item \textbf{MCP tools:} External MCP servers expose additional tools, accessible as \texttt{/mcp\_\_servername\_\_toolname}.
  \end{itemize}

\begin{lstlisting}
# Subagent with restricted tools:
---
name: review
allowed-tools: Read, Grep, Glob
---
\end{lstlisting}
\end{frame}


%%%%%%%%%%%%%%%%%%%%%%%%%%%%%%%%%%%%%%%%%%%%%%%%%%%%%%%%%%%%%%%%%%%%%%%%%%%%%%%%%%
\begin{frame}[fragile]\frametitle{Skills: Reusable Behavior Packages}
  \begin{itemize}
    \item \textbf{What:} Markdown instruction folders loaded on-demand by Claude when the task context matches. Unlike \texttt{CLAUDE.md} (always loaded), skills are lazy: Claude picks them only when relevant.
    \item \textbf{Skills ship with Claude Code:} Anthropic publishes official skills for docx, pdf, pptx, xlsx creation at \texttt{github.com/anthropics/claude-code} --- usable in any project.
    \item \textbf{Discovery locations} (all checked at startup):
      \begin{itemize}
        \item \texttt{.claude/skills/<n>/SKILL.md}: project-local
        \item \texttt{\textasciitilde/.claude/skills/<n>/SKILL.md}: global
      \end{itemize}
    \item \textbf{Activation:} Claude reads only the \texttt{description} field at startup. If the task matches, Claude loads the full SKILL.md. No manual invocation needed.
  \end{itemize}

\begin{lstlisting}
# .claude/skills/pytest-patterns/SKILL.md
---
name: pytest-patterns
description: >
  Best practices for writing pytest tests for this
  project. Use when generating or reviewing Python
  unit tests for mychatbot.
---
## Patterns
- Use @pytest.fixture for shared setup
- Mock external API calls with pytest-mock
- Use Faker for synthetic test data
\end{lstlisting}
\end{frame}


%%%%%%%%%%%%%%%%%%%%%%%%%%%%%%%%%%%%%%%%%%%%%%%%%%%%%%%%%%%%%%%%%%%%%%%%%%%%%%%%%%
\begin{frame}[fragile]\frametitle{Commands: Slash-Invokable Shortcuts}
  \begin{itemize}
    \item \textbf{What:} Reusable prompt templates invoked with \texttt{/} in the REPL. Eliminate repetitive prompt writing for common workflows.
    \item \textbf{Built-in commands:}
      \begin{itemize}
        \item \texttt{/init}: scan project, generate \texttt{CLAUDE.md}
        \item \texttt{/model}: switch active model
        \item \texttt{/compact}: summarize context to save tokens
        \item \texttt{/cost}: show current session token usage and cost
        \item \texttt{/help}: list all available commands (including custom)
        \item \texttt{/hooks}: list all configured hooks
        \item \texttt{/memory}: view or clear auto-generated memory
      \end{itemize}
    \item \textbf{Create custom commands} as \texttt{.claude/commands/<n>.md} (project) or \texttt{\textasciitilde/.claude/commands/<n>.md} (global).
    \item \textbf{\texttt{\$ARGUMENTS}:} Anything typed after \texttt{/command-name} is injected here: makes commands parameterized.
    \item \textbf{Team sharing:} Commit \texttt{.claude/commands/} to Git: the whole team gets it.
  \end{itemize}

\begin{lstlisting}
# .claude/commands/test-gen.md
---
description: Generate pytest tests for a source file
allowed-tools: Read, Edit, Write, Bash
---
Generate comprehensive pytest tests for $ARGUMENTS.
Read the pytest-patterns skill first, then write tests
covering happy path, edge cases, and error cases.
\end{lstlisting}
\end{frame}

%%%%%%%%%%%%%%%%%%%%%%%%%%%%%%%%%%%%%%%%%%%%%%%%%%%%%%%%%%%%%%%%%%%%%%%%%%%%%%%%%%
\begin{frame}[fragile]\frametitle{Permissions: Granular Action Control}
  \begin{itemize}
    \item \textbf{What:} Per-tool access policy: \texttt{allow} (auto-run), \texttt{ask} (prompt user), or \texttt{deny} (block). Defaults: most tools \texttt{allow}; \texttt{.env} files and \texttt{secrets/} denied by default.
    \item \textbf{Scope:} Set in \texttt{.claude/settings.json} (project) or \texttt{\textasciitilde/.claude/settings.json} (user). Subagent \texttt{allowed-tools} overrides project settings for that agent.
    \item \textbf{Enterprise managed:} \texttt{/etc/claude-code/managed-settings.json} --- cannot be bypassed by any user or project config.
    \item \textbf{Wildcard patterns:} \texttt{Bash(git:*)} allows all git commands; \texttt{Edit(src/**)} allows edits only in \texttt{src/}.
    \item \textbf{Caution: \texttt{--dangerously-skip-permissions}} bypasses the entire permission system. Never use in shared, CI, or production environments.
  \end{itemize}

\begin{lstlisting}
{ "permissions": {
    "allow": [
      "Read", "Glob", "Grep",
      "Bash(npm run:*)",
      "Bash(git commit:*)",
      "Edit(src/**)"
    ],
    "deny": [
      "Read(.env*)",
      "Read(secrets/**)",
      "Bash(rm -rf:*)",
      "Bash(sudo:*)"
    ],
    "ask": [
      "WebFetch",
      "Bash(docker:*)"
    ]
  }
}
\end{lstlisting}
\end{frame}


%%%%%%%%%%%%%%%%%%%%%%%%%%%%%%%%%%%%%%%%%%%%%%%%%%%%%%%%%%%%%%%%%%%%%%%%%%%%%%%%%%
\begin{frame}[fragile]\frametitle{Hooks: Deterministic Automation}
  \begin{itemize}
    \item \textbf{What:} Shell/prompt/agent handlers that execute automatically at specific lifecycle events. Unlike prompts, hooks \textit{always} run --- regardless of model behaviour.
    \item \textbf{Why:} Memory is probabilistic. Tooling is deterministic. Use hooks for anything that must always happen: auto-formatting, security checks, test runs.
    \item \textbf{9 lifecycle events:} PreToolUse, PostToolUse, Stop, SubagentStop, SessionStart, SessionEnd, UserPromptSubmit, Notification, PreCompact.
    \item \textbf{Three handler types:} \texttt{command} (shell script), \texttt{prompt} (Claude model evaluation), \texttt{agent} (subagent with tool access).
    \item \textbf{Exit codes:} 0 = proceed, 2 = block + send reason to Claude.
  \end{itemize}

\begin{lstlisting}
{ "hooks": {
    "PostToolUse": [
      { "matcher": "Edit|Write",
        "hooks": [{ "type": "command",
          "command": "black \"$CLAUDE_FILE_PATHS\" 2>/dev/null || true"
        }]
      }
    ],
    "PreToolUse": [
      { "matcher": "Bash",
        "hooks": [{ "type": "command",
          "command": "INPUT=$(cat); echo $INPUT | jq -r '.tool_input.command' | grep -qE 'rm -rf' && echo 'Blocked' >&2 && exit 2 || exit 0"
        }]
      }
    ]
  }
}
\end{lstlisting}
\end{frame}


%%%%%%%%%%%%%%%%%%%%%%%%%%%%%%%%%%%%%%%%%%%%%%%%%%%%%%%%%%%%%%%%%%%%%%%%%%%%%%%%%%
\begin{frame}[fragile]\frametitle{MCP Servers: External Tool Integration}
  \begin{itemize}
    \item \textbf{What:} Model Context Protocol servers expose external tools to Claude Code. Any MCP-compatible server becomes a tool Claude can call natively.
    \item \textbf{Why:} Extends Claude Code beyond the filesystem: databases, APIs, cloud services, monitoring tools, internal systems.
    \item \textbf{Popular MCP servers:} GitHub, Slack, PostgreSQL, Playwright, Brave Search, Linear, Jira, Notion, Sentry, AWS, Docker.
    \item \textbf{First-class CLI:} \texttt{claude mcp add}, \texttt{claude mcp list}, \texttt{claude mcp remove} --- no manual JSON editing needed.
    \item \textbf{Scope flags:} \texttt{--scope local} (default, personal), \texttt{--scope project} (shared via \texttt{.mcp.json}), \texttt{--scope user} (all your projects).
    \item \textbf{Tool naming:} MCP tools appear as \texttt{/mcp\_\_servername\_\_toolname} slash commands. Use \texttt{mcp\_\_github} in permissions to control all tools from one server.
    \item \textbf{Build your own:} Any HTTP/stdio process implementing the MCP spec; SDKs in TypeScript, Python, Go. Use FastMCP for Python servers.
  \end{itemize}

\end{frame}

%%%%%%%%%%%%%%%%%%%%%%%%%%%%%%%%%%%%%%%%%%%%%%%%%%%%%%%%%%%%%%%%%%%%%%%%%%%%%%%%%%
\begin{frame}[fragile]\frametitle{MCP Servers: Configuration}

\begin{lstlisting}
# Add via CLI (recommended):
claude mcp add github \
  --transport http \
  https://api.githubcopilot.com/mcp \
  -H "Authorization: Bearer $GITHUB_TOKEN"

# Or edit .mcp.json directly (committed to repo):
{ "mcpServers": {
    "github": {
      "command": "npx",
      "args": ["-y","@modelcontextprotocol/server-github"],
      "env": { "GITHUB_TOKEN": "$GITHUB_TOKEN" }
    },
    "postgres": {
      "command": "npx",
      "args": ["-y","@benborla29/mcp-server-postgresql",
               "postgresql://localhost/mydb"]
    },
    "sequential-thinking": {
      "type": "stdio",
      "command": "npx",
      "args":["-y",
        "@modelcontextprotocol/server-sequential-thinking"]
    }
  }
}
\end{lstlisting}
\end{frame}


%%%%%%%%%%%%%%%%%%%%%%%%%%%%%%%%%%%%%%%%%%%%%%%%%%%%%%%%%%%%%%%%%%%%%%%%%%%%%%%%%%
\begin{frame}[fragile]\frametitle{Plugins: Distributable Workflow Bundles}
  \begin{itemize}
    \item \textbf{What:} Installable packages that combine skills, subagents, hooks, commands, and MCP configs. Solve the ``works-on-my-machine'' problem for team tooling.
    \item \textbf{Why:} Instead of every team member manually configuring their \texttt{.claude/} folder, package everything once and share.
    \item \textbf{Plugin structure:} A folder containing \texttt{.claude-plugin/plugin.json} + optional \texttt{commands/}, \texttt{agents/}, \texttt{skills/}, \texttt{hooks/}, \texttt{.mcp.json}, \texttt{scripts/}.
    \item \textbf{Namespacing:} Skills in plugins are namespaced \texttt{/pluginName:skillName} to avoid collisions.
    \item \textbf{Install / share:}
      \begin{itemize}
        \item Local: \texttt{claude plugin add --path ./my-plugin}
        \item Marketplace: \texttt{claude plugin install <name>}
        \item Auto-install for team: list in \texttt{.claude/settings.json}
      \end{itemize}
    \item \textbf{Community marketplace:} Growing library of open-source plugins at \href{https://github.com/hesreallyhim/awesome-claude-code}{\texttt{github.com/hesreallyhim/awesome-claude-code}}.
  \end{itemize}

\begin{lstlisting}
# .claude-plugin/plugin.json
{ "name": "my-team-tools",
  "version": "1.0.0",
  "description": "Shared dev workflow for our team",
  "author": { "name": "Engineering Team" } }
\end{lstlisting}
\end{frame}


%%%%%%%%%%%%%%%%%%%%%%%%%%%%%%%%%%%%%%%%%%%%%%%%%%%%%%%%%%%%%%%%%%%%%%%%%%%%%%%%%%
\begin{frame}[fragile]\frametitle{Models: Tier \& Configuration}
  \begin{itemize}
    \item \textbf{Three tiers}, each with a different cost-quality tradeoff:
      \begin{itemize}
        \item \textbf{Haiku 4.5}: fastest, cheapest; ideal for subagents doing exploration and simple edits.
        \item \textbf{Sonnet 4.6}: balanced default; covers most daily coding and refactoring tasks.
        \item \textbf{Opus 4.6}: most capable; use for architecture design, complex reasoning, security review.
      \end{itemize}
    \item Choose models with \textbf{long context window} and \textbf{tool use (function calling)} support.
    \item \textbf{Switch model} in REPL: \texttt{/model claude-opus-4-6}. Or \texttt{--model} CLI flag at launch.
    \item \textbf{Configure in \texttt{.claude/settings.json}:} \texttt{"model": "claude-sonnet-4-6"}.
    \item \textbf{Per-subagent model:} Each custom subagent can specify its own model tier --- e.g., \texttt{debug} uses Haiku for fast file search; \texttt{review} uses Opus for thoroughness.
    \item \textbf{Extended thinking:} Available on Sonnet and Opus; toggle with \texttt{Alt+T} or include in subagent definition.
    \item \textbf{Context window:} 200K tokens standard; up to \textbf{1M tokens} with Opus 4.6+ (beta).
  \end{itemize}

\end{frame}

%%%%%%%%%%%%%%%%%%%%%%%%%%%%%%%%%%%%%%%%%%%%%%%%%%%%%%%%%%%%%%%%%%%%%%%%%%%%%%%%%
\begin{frame}[fragile]\frametitle{Models: Configuration Example}

\begin{lstlisting}
# .claude/settings.json
{
  "$schema":
   "https://json.schemastore.org/claude-code-settings.json",
  "model": "claude-sonnet-4-6",
  "permissions": {
    "allow": ["Read", "Bash(npm run:*)", "Edit(src/**)"],
    "deny":  ["Read(.env*)", "Bash(rm -rf:*)"]
  },
  "env": {
    "NODE_ENV": "development"
  }
}

# Per-subagent model override in .claude/agents/review.md:
---
name: review
model: claude-opus-4-6
allowed-tools: Read, Grep, Glob
---

# Switch model in session:
/model claude-opus-4-6

# Launch with specific model:
claude --model claude-haiku-4-5
\end{lstlisting}
\end{frame}
